% الجزء: sets-functions-relations
% الفصل: sets
% القسم: unions-and-intersections

\documentclass[../../../include/open-logic-section]{subfiles}

\begin{document}

\olfileid[ar]{sfr}{set}{uni}
\olsection{الاتحادات والتقاطعات}

\begin{explain}
قدّمنا في \olref[sfr][set][bas]{sec} تعريفات للمجموعات بالتجريد،
أي تعريفات من الصورة $\Setabs{x}{\phi(x)}$. ونستعمل هنا
خاصية ما~$\phi$، وقد تشير هذه الخاصية إلى مجموعات سبق أن عرّفناها.
فمثلًا، إذا كانت $A$ و~$B$ مجموعتين، فإن المجموعة
$\Setabs{x}{x \in A \lor x \in B}$ تتألف من جميع الكائنات التي تُعَدّ من
!!{element}s $A$ أو~$B$، أي إنها المجموعة التي تجمع !!{element}s
$A$ و~$B$. ويمكن تمثيل ذلك بصريًا كما في \olref{fig:union}، حيث تدل
المنطقة المظللة على !!{element}s المجموعتين $A$ و~$B$ مجتمعةً.

\begin{figure}
  \olasset{assets/diagrams/union.tikz}
  \caption{اتحاد المجموعتين $A \cup B$ هو المجموعة المكوّنة من !!{element}s
   $A$ و~$B$ معًا.}
  \ollabel{fig:union}
\end{figure}

وهذه العملية على المجموعات---أي جمعها---شائعة ومفيدة جدًا،
ولذلك نخصها باسم اصطلاحي ورمز.
\end{explain}

\begin{defn}[الاتحاد]
\emph{اتحاد} مجموعتين $A$ و$B$، ويكتب $A \cup B$، هو مجموعة
جميع الأشياء التي تكون من !!{element}s $A$ أو $B$ أو كلتيهما.
\[
A \cup B = \Setabs{x}{x \in A \lor x \in B}
\]
\end{defn}

\begin{ex}
بما أن تكرار !!{element}s لا يعتد به، فإن اتحاد مجموعتين تشتركان في
!!a{element} لا يحتوي ذلك !!{element} إلا مرة واحدة؛ فمثلًا،
$\{ a, b, c\} \cup \{ a, 0, 1\} = \{a, b, c, 0, 1\}$.

اتحاد مجموعة وإحدى مجموعاتها الجزئية هو المجموعة الأكبر نفسها: $\{a,
b, c \} \cup \{a \} = \{a, b, c\}$.

اتحاد مجموعة والمجموعة الخالية هو المجموعة نفسها: $\{a,
b, c \} \cup \emptyset = \{a, b, c \}$.
\end{ex}

\begin{prob}
أثبت أنه إذا كان $A \subseteq B$، فإن $A \cup B = B$.
\end{prob}

\begin{explain}
يمكننا أيضًا النظر في عملية «مزدوجة» للاتحاد. وهي العملية التي تكوّن
مجموعة جميع !!{element}s التي تُعَدّ من !!{element}s~$A$ ومن !!{element}s~$B$
أيضًا. وتسمى هذه العملية \emph{التقاطع}، ويمكن تمثيلها كما في
\olref{fig:intersection}.
\begin{figure}
  \olasset{assets/diagrams/intersection.tikz}
  \caption{تقاطع المجموعتين $A \cap B$ هو المجموعة المكوّنة من
    !!{element}s مشتركة بينهما.}
  \ollabel{fig:intersection}
\end{figure}
\end{explain}

\begin{defn}[التقاطع]
\emph{تقاطع} مجموعتين $A$ و$B$، ويكتب $A \cap B$، هو مجموعة
جميع الأشياء التي تكون من !!{element}s كل من $A$ و~$B$.
\[
A \cap B = \Setabs{x}{x \in A \land x \in B}
\]
وتُسمّى مجموعتان \emph{متباينتين} إذا كان تقاطعهما
خاليًا. وهذا يعني أنه لا توجد بينهما !!{element}s مشتركة.
\end{defn}

\begin{ex}
إذا لم تكن بين مجموعتين !!{element}s مشتركة، كان تقاطعهما خاليًا:
$\{ a, b, c\} \cap \{ 0, 1\} = \emptyset$.

وإذا كانت بين مجموعتين !!{element}s مشتركة، كان تقاطعهما مجموعة
تلك العناصر كلها: $\{a, b, c \} \cap \{a, b, d \} = \{a, b\}$.

تقاطع مجموعة وإحدى مجموعاتها الجزئية هو المجموعة الأصغر نفسها:
$\{a, b, c\} \cap \{a, b\} = \{a, b\}$.

تقاطع أي مجموعة والمجموعة الخالية خالٍ: $\{a, b, c \}
\cap \emptyset = \emptyset$.
\end{ex}

\begin{prob}
أثبت إثباتًا صارمًا أنه إذا كان $A \subseteq B$، فإن
$A \cap B = A$.
\end{prob}

\begin{explain}
كما يمكننا تكوين اتحاد أو تقاطع لأكثر من مجموعتين. ومن الطرق الأنيقة
لمعالجة ذلك بوجه عام ما يأتي: افترض أنك تجمع كل المجموعات التي تريد
تكوين اتحادها (أو تقاطعها) في مجموعة واحدة. عندئذ يمكننا تعريف اتحاد
مجموعاتنا الأصلية كلها بأنه مجموعة الكائنات التي تنتمي إلى !!{element}
واحد على الأقل من تلك المجموعة، وتعريف التقاطع بأنه مجموعة كل الكائنات
التي تنتمي إلى كل !!{element} من تلك المجموعة.
\end{explain}

\begin{defn}
إذا كانت $A$ مجموعة من المجموعات، فإن $\bigcup A$ هي المجموعة التي
تضم !!{element}s جميع !!{element}s~$A$:
\begin{align*}
\bigcup A & = \Setabs{x}{x \text{ ينتمي إلى !!a{element} من } A},
\text{ أي}\\
& = \Setabs{x}{\text{يوجد } B \in A
  \text{ بحيث } x \in B}
\end{align*}
\end{defn}

\begin{defn}
إذا كانت $A$ مجموعة من المجموعات، فإن $\bigcap A$ هي مجموعة الكائنات
المشتركة بين جميع عناصر~$A$:
\begin{align*}
\bigcap A & = \Setabs{x}{x \text{ ينتمي إلى كل !!{element} من } A},
\text{ أي}\\
 & = \Setabs{x}{\text{لكل } B \in A, x \in B}
\end{align*}
\end{defn}

\begin{ex}
افترض أن $A = \{ \{ a, b \}, \{ a, d, e \}, \{ a, d \} \}$.
عندئذ $\bigcup A = \{ a, b, d, e \}$ و$\bigcap A = \{ a \}$.
\end{ex}
\begin{prob}
	بيّن أنه إذا كانت $A$ مجموعة وكان $A \in B$، فإن
$A \subseteq \bigcup B$.
\end{prob}

ويمكننا فعل الأمر نفسه لمتتالية من المجموعات $A_1$، $A_2$، \dots
\begin{align*}
\bigcup_i A_i & = \Setabs{x}{x \text{ ينتمي إلى إحدى المجموعات } A_i}\\
\bigcap_i A_i & = \Setabs{x}{x \text{ ينتمي إلى كل } A_i}.
\end{align*}

وعندما تكون لدينا \emph{مجموعة فهرسة} للمجموعات، أي مجموعة ما $I$
بحيث ننظر في $A_i$ لكل $i \in I$، يمكننا أيضًا استعمال الاختصارين
الآتيين:
\begin{align*}
	\bigcup_{i \in I} A_i & = \bigcup \Setabs{A_i }{i \in I}\\
	\bigcap_{i \in I} A_i & = \bigcap\Setabs{A_i}{i \in I}
\end{align*}

وأخيرًا، قد نريد النظر في مجموعة جميع !!{element}s الموجودة في~$A$
وغير الموجودة في~$B$. ويمكن تمثيل ذلك كما في \olref{difference}.

\begin{figure}
  \olasset{assets/diagrams/difference.tikz}
  \caption{فرق المجموعتين $A \setminus B$ هو المجموعة المكوّنة من
    !!{element}s~$A$ التي لا تُعَدّ أيضًا من !!{element}s~$B$.}
  \ollabel{difference}
\end{figure}

\begin{defn}[الفرق]
\emph{فرق المجموعتين}~$A \setminus B$ هو مجموعة جميع !!{element}s
$A$ التي لا تُعَدّ أيضًا من !!{element}s~$B$، أي
\[
A\setminus B = \Setabs{x}{x\in A \text{ و } x \notin B}.
\]
\end{defn}

\begin{prob}
	أثبت أنه إذا كان $A \subsetneq B$، فإن $B \setminus A \neq \emptyset$.
\end{prob}

\end{document}
